\newpage
\section{Como correr o projeto (Docker -- Deploy em Produção)}

O ecossistema foi desenhado para ser agnóstico ao ambiente. Através do GitHub Container Registry (GHCR), o \textit{deploy} é efetuado de forma imutável e automática.

\subsection{Comando de Deploy}
Este comando descarrega o manifesto de orquestração e levanta todos os serviços (Gateway, Microserviços, Base de Dados e Frontend) sem necessidade de dependências locais.



\textit{Windows (PowerShell):}
\begin{codeblock}
curl.exe -L -o docker-compose.deploy.yml "https://raw.githubusercontent.com/kanekiTakitos/Nexus-Estates/main/infrastructure/docker-compose.deploy.yml" ; docker compose -f docker-compose.deploy.yml up -d --pull always
\end{codeblock}




\textit{Linux / Mac / WSL:}
\begin{codeblock}
curl -L -o docker-compose.deploy.yml "https://raw.githubusercontent.com/kanekiTakitos/Nexus-Estates/main/infrastructure/docker-compose.deploy.yml" && docker compose -f docker-compose.deploy.yml up -d --pull always
\end{codeblock}


\subsection{Acessos e Credenciais}
Após o arranque, os serviços estarão operacionais nos seguintes \textit{endpoints}:

\begin{itemize}[label=\textbullet, noitemsep]
    \item \textbf{Frontend:} \url{http://localhost:3000}
    \item \textbf{API Gateway:} \url{http://localhost:8080}
    \item \textbf{Swagger UI:} \url{http://localhost:8080/swagger-ui.html}
\end{itemize}

\textbf{Credenciais de Teste:} \\
\texttt{Email: dev@nexus.com} \\
\textbf{Password:} \texttt{12345}

\subsection{Gestão do Ciclo de Vida}

\textbf{Parar e limpar dados (Reset):}
\begin{codeblock}
docker compose -f docker-compose.deploy.yml down -v
\end{codeblock}

\textbf{Remoção Completa (Cleanup Total):}
Remove contentores, volumes, imagens e o ficheiro YAML gerado pelo \textit{curl}.

\textit{Linux / Mac / WSL:}
\begin{codeblock}
docker compose -f docker-compose.deploy.yml down -v --rmi all &&   rm docker-compose.deploy.yml
\end{codeblock}

\textit{Windows (PowerShell):}
\begin{codeblock}
docker compose -f docker-compose.deploy.yml down -v --rmi all ; del docker-compose.deploy.yml
\end{codeblock}

\subsection{Arquitetura de Rede e Segurança}
A infraestrutura utiliza uma rede isolada (\texttt{nexus-net}). 
\begin{itemize}[label=\textbullet, noitemsep]
    \item \textbf{Isolamento:} Apenas o Frontend e o Gateway expõem portas ao \textit{host}.
    \item \textbf{Service Discovery:} A comunicação \textit{inter-service} é feita via DNS interno do Docker.
    \item \textbf{Segurança:} Bases de dados e \textit{brokers} de mensagens estão inacessíveis fora da rede interna.
\end{itemize}

\section{Ambiente de Desenvolvimento (DEV)}
Para desenvolvimento ativo com \textit{Hot Reload} e \textit{debug}:

\subsection{Infraestrutura Base}
Levanta apenas os serviços de suporte (Base de Dados e RabbitMQ).
\begin{codeblock}
docker compose -f infrastructure/docker-compose.yml up -d
\end{codeblock}

\subsection{Backend e Frontend}
Comandos para execução local dos serviços em modo de desenvolvimento.
\begin{codeblock}
# Terminal 1 (Java/Spring)
mvn -pl api-gateway -am spring-boot:run

# Terminal 2 (Next.js)
cd frontend && bun dev
\end{codeblock}

Atualmente usamos o IDE de jetbrains que facilita a instalação e compilação.